% !TEX root = ../main.tex
% T-ChatKBQA System Architecture Diagram (TikZ)
% Included by sections/01_chatkbqa.tex

\begin{tikzpicture}[
    node distance=0.18cm,
    % === Module boxes (main pipeline) ===
    module/.style={
        draw=blue!60!black, fill=blue!8, rounded corners=3pt,
        text width=1.85cm, minimum height=1.0cm, align=center, font=\tiny\bfseries, inner sep=2pt
    },
    submodule/.style={
        draw=blue!40, fill=blue!4, rounded corners=2pt,
        text width=1.9cm, minimum height=0.5cm, align=center, font=\fontsize{5}{6}\selectfont, inner sep=1pt
    },
    % === Training pipeline boxes ===
    trainmod/.style={
        draw=teal!70!black, fill=teal!8, rounded corners=2pt,
        text width=1.45cm, minimum height=0.55cm, align=center, font=\fontsize{5}{6}\selectfont, inner sep=2pt
    },
    % === Arrows ===
    arrow/.style={-stealth, blue!60!black, thick},
    arrowgray/.style={-stealth, gray, thin},
    arrowteal/.style={-stealth, teal!60!black, thick},
    % === Labels on arrows ===
    alabel/.style={font=\fontsize{5}{6}\selectfont\itshape, fill=white, inner sep=0.5pt},
    % === Group backgrounds ===
    group/.style={draw=blue!25, rounded corners=5pt, inner sep=4pt},
    traingroup/.style={draw=teal!50!black, rounded corners=5pt, inner sep=4pt, fill=teal!3},
]

% ============================================================
% TRAINING DATA PIPELINE (bottom)
% ============================================================
\begin{scope}[local bounding box=trainbox]
    \node[trainmod] (T0) at (0,0) {Zenodo\\FB+CVT+REV\\$\sim$244M trips};
    \node[trainmod, right=0.15cm of T0] (T1) {Fact\\Mining\\(22 rels)};
    \node[trainmod, right=0.15cm of T1] (T2) {Template\\Bank\\(15 tmpl)};
    \node[trainmod, right=0.15cm of T2] (T3) {Dist.\\Filter\\(2,061)};
    \node[trainmod, right=0.15cm of T3] (T4) {Train\\Export\\(SFT fmt)};

    \draw[arrowteal] (T0) -- (T1);
    \draw[arrowteal] (T1) -- (T2);
    \draw[arrowteal] (T2) -- (T3);
    \draw[arrowteal] (T3) -- (T4);
\end{scope}

\begin{pgfonlayer}{background}
    \node[traingroup, fit=(T0)(T1)(T2)(T3)(T4),
          label={[teal!60!black, font=\tiny\bfseries]above:Training Data Construction}] {};
\end{pgfonlayer}

% ============================================================
% MAIN INFERENCE PIPELINE
% ============================================================
\begin{scope}[local bounding box=mainpipe]
    % --- Input ---
    \node[module, fill=orange!15, draw=orange!70!black] (M0) at (0, 2.5)
        {NL Temporal\\Question\\\texttt{"Who was US}\\ \texttt{president..."}};

    % --- Agent Router ---
    \node[submodule, fill=yellow!15, draw=yellow!60!black, below=0.15cm of M0] (AG)
        {Agent Router\\(before/after/first/last)};

    % --- Module 2 ---
    \node[module, right=0.15cm of M0] (M1)
        {M2: Logical\\Form Gen.\\LLaMA-2-7B\\+ LoRA (k=5)};

    % --- Module 3 ---
    \node[module, right=0.15cm of M1] (M2)
        {M3: Extended\\Retrieval\\Entity + Rel\\+ Timestamp};

    % --- Module 4 ---
    \node[module, right=0.15cm of M2] (M3)
        {M4: Temporal\\SPARQL Exec\\Virtuoso/\\Freebase};

    % --- Output ---
    \node[module, fill=green!15, draw=green!60!black, right=0.15cm of M3] (M4)
        {Answer\\+ Provenance\\+ Interpr.\\SPARQL};
\end{scope}

% Main pipeline arrows
\draw[arrow] (M0) -- (M1) node[alabel, midway, above] {\tiny signal};
\draw[arrow] (M1) -- (M2) node[alabel, midway, above] {\tiny S-expr};
\draw[arrow] (M2) -- (M3) node[alabel, midway, above] {\tiny res. IDs};
\draw[arrow] (M3) -- (M4) node[alabel, midway, above] {\tiny results};

% Agent → routing arrows
\draw[arrow, shorten >=2pt] (M0.south) -- (AG.north);
\draw[arrowgray, shorten >=2pt] (AG.east) to[out=0, in=270] (M1.south west);

% Training → LLM arrow
\draw[arrowteal, dashed] (T4.north) to[out=90, in=270] (M1.south)
    node[alabel, midway, right, fill=white] {\tiny LoRA};

% ============================================================
% DETAILED SUB-BOXES for Module 3
% ============================================================
\begin{scope}[local bounding box=retdetails]
    \node[submodule, below=0.2cm of M2, text width=1.7cm] (R1)
        {Entity\\ELQ/SimCSE};
    \node[submodule, right=0.15cm of R1, text width=1.7cm] (R2)
        {Relation\\SimCSE};
    \node[submodule, right=0.15cm of R2, text width=1.7cm] (R3)
        {Timestamp\\Date Norm.};
\end{scope}
\draw[arrowgray, shorten >=2pt] (M2.south) -- (R1.north);
\draw[arrowgray, shorten >=2pt] (M2.south) -- (R2.north);
\draw[arrowgray, shorten >=2pt] (M2.south) -- (R3.north);

% ============================================================
% TEMPORAL OPERATORS LEGEND
% ============================================================
\node[draw=purple!60!black, fill=purple!5, rounded corners=3pt,
      inner sep=3pt, font=\fontsize{5}{6}\selectfont, text width=4.2cm, align=left,
      below=0.1cm of trainbox.south east, anchor=north east] (OP)
    {\textbf{Operators:}
     \texttt{TC}(range) $|$ \texttt{ARGMAX}(last) $|$ \texttt{ARGMIN}(first) $|$
     \texttt{gt/ge/lt/le}(cmp) $|$ \texttt{JOIN}};

% ============================================================
% GROUP BACKGROUNDS
% ============================================================
\begin{pgfonlayer}{background}
    \node[group, fit=(M0)(AG)(M1)(M2)(M3)(M4),
          label={[blue!60!black, font=\small\bfseries]above:T-ChatKBQA System Architecture}] {};
\end{pgfonlayer}

\end{tikzpicture}
